% LaTeX header-includes for PDF generation (see pdf_export.py). Requires
% xelatex (not pdflatex) since it uses fontspec to pull in real installed
% system fonts by name rather than LaTeX's built-in Computer Modern.
%
% Body text: Open Sans. Headings (all of #/##/###/#### in the generated
% Markdown, i.e. \section/\subsection/\subsubsection/\paragraph once Pandoc
% converts them): Lora Bold. Table column headings: Open Sans SemiBold.
\usepackage{fontspec}
\setmainfont{Open Sans}
\newfontfamily\headingfont{Lora}
\newfontfamily\tableheadfont{Open Sans SemiBold}

\usepackage{titlesec}
\titleformat{\section}{\headingfont\bfseries\Large}{}{0pt}{}
\titleformat{\subsection}{\headingfont\bfseries\large}{}{0pt}{}
\titleformat{\subsubsection}{\headingfont\bfseries\normalsize}{}{0pt}{}
% \paragraph defaults to a run-in heading (same line as the text that
% follows) — [block] makes it sit on its own line like the other levels,
% since #### marks a distinct finding heading in the report, not a run-in
% label.
\titleformat{\paragraph}[block]{\headingfont\bfseries\normalsize}{}{0pt}{}
\titlespacing*{\paragraph}{0pt}{2ex}{1ex}

% Table column headings (Open Sans SemiBold, via \tableheadfont above) are
% applied per-cell by pandoc_table_header_font.lua, not here — redefining
% \toprule/\midrule to inject a font switch broke booktabs/longtable's
% internal \noalign handling ("Misplaced \noalign", confirmed empirically).
% The Lua filter operates on the Pandoc AST before LaTeX serialization, so
% it only ever touches header-cell text, never the surrounding table
% structure commands.
